% Pursuing Stacks (cotes 134-2 to 134-8): archive cote 134-2, site page 59 (PDF page 60).
% appendix (letters to Breen). Draft compared: ps1.tex:1423-1438.
% AI-transcribed by Claude Opus 5.5 (claude-opus-5-5, adaptive thinking, effort high), Message Batches API, batch ps1 (msgbatch_01Rv5hCaNTtUt4JHx1yqoHPL), 2026-11-12; not yet checked by eye.
\scanpage{134-2/59}
\mspage{-\hand{\unsure{1}}4 -}

coefficients - Verdier, \emph{exposé} Bourbaki (SNLM 300) - and for the étale
cohomology of schemes for ``discrete'' coefficients (\deleted{with} ``$\ell$-adic'' \inserted{or} torsion)
prime to the residual characteristic\hand{\unsure{s}}\over{-}{(} SGA5 \over{-}{)}, finally, for coherent sheaves
on analytic spaces \over{-}{(} Verdier-Ruget\over{-}{)}\hand{.} (The formalism remains to be develop\over{pe}{ed}
in the crystalline context, and in characteristic $0$ in the context of
\inserted{stratified} module\hand{s} \deleted{stratifications and} \inserted{with} singularities, à la Déligne, with perhaps \deleted{-}
\hand{-} over the field $\mathbb{C}$ - the introduction of additional Hodge structures, finally
in the context of moti\over{fs}{ves}\over{.}{;} I am \deleted{persuaded} \inserted{convinced} that it exists \deleted{in some other places} \inserted{about anywhere}
- \deleted{in all the places, perhaps where} \inserted{maybe, wherever} there is a formalism of a cohomological
nature.)

Working in étale cohomology on a separated scheme of finite type \over{on}{over} a
field $k$ , say, with a ring of coefficients $\Lambda$ of torsion prime to the
characteristic, the complex of sheaves $f^{!}(\Lambda_e)$ (where $\hat{f} : X \to \mathrm{Spec}\ k = e$)
plays the rôle of \emph{complex of singular chains on} $X$ \emph{with coefficients in} $\Lambda$ ,
and $Rf_{!}$ $(f^{!}\Lambda_e)$ plays the rôle of a \emph{homology} $H_{*}(X/e)$ , \emph{vis a vis} of course,
of coefficients on $e$ which are complexes of $\Lambda$-modules. You can easily
justify this assertion with the help of the ``global duality theorems'', by
one or two tricks which I spare you here.

\uline{REMARKS} (1) There is no \inserted{need to} truncat\over{\unsure{ion}}{ion}, it works in all dimensions.
(2) This is related (at least \deleted{from a} \inserted{as far as the} philosoph\over{ical}{y} \deleted{point of view} \inserted{is concerned}) to the
fact that for the various types of coefficients (under conditions of
``constructibility'') one has a theorem of ``biduality'', at least if one allows
resolution of singularities (but Deligne has told me I believe that he knows
a proof without that)\hand{,} with values in a ``duali\over{\unsure{s}}{\unsure{z}}ing complex'' $K_e$ (on $e$),
$K_X$ (on $X$) . If for example $\Lambda$ is ``self-dualising'' (or Gorenstein) for
example $\Lambda = \mathbb{Z}/n\mathbb{Z}$\deleted{\illegible{1}}\hand{,} one can take $K_e = \Lambda$ , therefore the dualising complex
$K_X = f^{!}(K_e)$ is nothing else than the ``complex of singular chains with
coefficients in $\Lambda$ ''.
(3) One can do the same thing for coefficients such as $\mathbb{Z}_{\hand{$\ell$}}$ (Jouanolou, \hand{thesis -}
non published \deleted{\unsure{these|thesis}}\hand{,} I fear!)\hand{.}

\mspage{\hand{59}}
